% Lambek--Scott 1980 L_1, section 1 pp.216--218; biconditionals split into rules.
\calculussheet{sheet:p10-internal-hol-ls1980}{Lambek--Scott 1980 product/power type theory}{Source: \cite[\S1, Definitions~1.1--1.2 and axioms, pp.~216--218]{LambekScott1980ITT}. Types are \(1,N,\Omega\), closed under \(A\times B\) and \(PA\); terms are exactly the ten clauses in Definition~\ref{def:p10-ls1980-language}. The explicitly nonarithmetic fragment omits \(N,0,S\) and Peano rows; the applied extension names additional interpreted object types and arrows as in source \S5 p.~242. Judgments \(p\vdash_Xq\) have formula free variables contained in the declared set \(X\); \(\vdash_Xp\) abbreviates \(\top\vdash_Xp\). All substitutions are typed and free-for/capture-avoiding. Write \(\neg p=(p\Rightarrow\bot)\), \(p\Leftrightarrow q=(p\Rightarrow q)\wedge(q\Rightarrow p)\), and \(a=_Ab\) for \(\forall u:PA\,(a\in u\Leftrightarrow b\in u)\), with \(u\) fresh. Equality is defined, not primitive; neither function/lambda types nor projections are primitive terms of pure \(\mathcal L_1\). Each source logical iff is displayed as its forward and reverse rules. Quantifier freshness is explicit; optional Peano rows require the NNO interpretation and the Boolean row requires Booleanness.}
\hbrule{rule:p10-internal-hol-ls1980-refl}{Reflexivity}{}{p\vdash_Xp}{\none}
\hbrule{rule:p10-internal-hol-ls1980-trans}{Transitivity}{p\vdash_Xq\qquad q\vdash_Xr}{p\vdash_Xr}{\none}
\hbrule{rule:p10-internal-hol-ls1980-weaken}{Variable weakening}{p\vdash_Xq}{p\vdash_Yq}{\(X\subseteq Y\); declared types are unchanged. No unused-variable deletion rule.}
\hbrule{rule:p10-internal-hol-ls1980-subst}{Typed substitution}{\phi(y)\vdash_{X\cup\{y:B\}}\psi(y)}{\phi(b)\vdash_X\psi(b)}{\(b:B\), \(\mathrm{FV}(b)\subseteq X\); \(b\) is free for \(y\) in both formulas. If \(y\in X\), its original assigned value is used when evaluating \(b\), and the premise receives the substituted value at \(y\).}
\hbrule{rule:p10-internal-hol-ls1980-top}{Logical (1)}{}{p\vdash_X\top}{\none}
\hbrule{rule:p10-internal-hol-ls1980-bottom}{Logical (2)}{}{\bot\vdash_Xp}{\none}
\hbrule{rule:p10-internal-hol-ls1980-meet-intro}{Logical (3), reverse}{p\vdash_Xq\qquad p\vdash_Xr}{p\vdash_Xq\wedge r}{\none}
\hbrule{rule:p10-internal-hol-ls1980-meet-left}{Logical (3), forward 1}{p\vdash_Xq\wedge r}{p\vdash_Xq}{\none}
\hbrule{rule:p10-internal-hol-ls1980-meet-right}{Logical (3), forward 2}{p\vdash_Xq\wedge r}{p\vdash_Xr}{\none}
\hbrule{rule:p10-internal-hol-ls1980-join-elim}{Logical (4), reverse}{p\vdash_Xr\qquad q\vdash_Xr}{p\vee q\vdash_Xr}{\none}
\hbrule{rule:p10-internal-hol-ls1980-join-left}{Logical (4), forward 1}{p\vee q\vdash_Xr}{p\vdash_Xr}{\none}
\hbrule{rule:p10-internal-hol-ls1980-join-right}{Logical (4), forward 2}{p\vee q\vdash_Xr}{q\vdash_Xr}{\none}
\hbrule{rule:p10-internal-hol-ls1980-residual-intro}{Logical (5), forward}{p\wedge q\vdash_Xr}{p\vdash_Xq\Rightarrow r}{\none}
\hbrule{rule:p10-internal-hol-ls1980-residual-elim}{Logical (5), reverse}{p\vdash_Xq\Rightarrow r}{p\wedge q\vdash_Xr}{\none}
\hbrule{rule:p10-internal-hol-ls1980-forall-intro}{Logical (6), reverse}{p\vdash_{X\cup\{y:B\}}\psi(y)}{p\vdash_X\forall y:B\,\psi(y)}{\(y\notin\mathrm{FV}(p)\); other free variables of \(\psi\) lie in \(X\). If \(y\in X\), its declared value is overridden when testing the bound variable.}
\hbrule{rule:p10-internal-hol-ls1980-forall-elim}{Logical (6), forward}{p\vdash_X\forall y:B\,\psi(y)}{p\vdash_{X\cup\{y:B\}}\psi(y)}{\(y\notin\mathrm{FV}(p)\); other free variables of \(\psi\) lie in \(X\).}
\hbrule{rule:p10-internal-hol-ls1980-exists-elim}{Logical (7), reverse}{\psi(y)\vdash_{X\cup\{y:B\}}p}{\exists y:B\,\psi(y)\vdash_Xp}{\(y\notin\mathrm{FV}(p)\); other free variables of \(\psi\) lie in \(X\).}
\hbrule{rule:p10-internal-hol-ls1980-exists-intro}{Logical (7), forward}{\exists y:B\,\psi(y)\vdash_Xp}{\psi(y)\vdash_{X\cup\{y:B\}}p}{\(y\notin\mathrm{FV}(p)\); other free variables of \(\psi\) lie in \(X\).}
\hbrule{rule:p10-internal-hol-ls1980-comprehension}{Nonlogical (a)}{}{\vdash_X\forall x:A\,(x\in\{x:A\mid\phi(x)\}\Leftrightarrow\phi(x))}{Parameters of \(\phi\) other than its displayed \(x\) lie in \(X\); binders are renamed to avoid capture, not to delete declared variables.}
\hbrule{rule:p10-internal-hol-ls1980-power-ext}{Nonlogical (b), power}{}{\vdash_{\varnothing}\forall u:PA\,\forall v:PA\,((\forall x:A\,(x\in u\Leftrightarrow x\in v))\Rightarrow u=_{PA}v)}{\(u,v,x\) are distinct bound variables.}
\hbrule{rule:p10-internal-hol-ls1980-proposition-ext}{Nonlogical (b), propositions}{}{\vdash_{\varnothing}\forall s:\Omega\,\forall t:\Omega\,((s\Leftrightarrow t)\Rightarrow s=_{\Omega}t)}{\(s,t\) are distinct bound variables.}
\hbrule{rule:p10-internal-hol-ls1980-unit}{Nonlogical (c), unit}{}{\vdash_{\varnothing}\forall x:1\,x=_1*}{\none}
\hbrule{rule:p10-internal-hol-ls1980-product-surj}{Nonlogical (c), decomposition}{}{\vdash_{\varnothing}\forall z:A\times B\,\exists x:A\,\exists y:B\,z=_{A\times B}\langle x,y\rangle}{\(z,x,y\) are distinct bound variables.}
\hbrule{rule:p10-internal-hol-ls1980-product-inj}{Nonlogical (c), injectivity}{}{\vdash_{\varnothing}\forall x:A\,\forall x':A\,\forall y:B\,\forall y':B\,(\langle x,y\rangle=_{A\times B}\langle x',y'\rangle\Rightarrow(x=_Ax'\wedge y=_By'))}{All four bound variables are distinct.}
\hbrule{rule:p10-internal-hol-ls1980-peano-zero}{Optional nonlogical (d), zero}{}{\vdash_{\varnothing}\forall x:N\,(Sx=_N0\Rightarrow\bot)}{Arithmetic extension only; \(N\) is interpreted as an NNO with its zero and successor.}
\hbrule{rule:p10-internal-hol-ls1980-peano-inj}{Optional nonlogical (d), successor}{}{\vdash_{\varnothing}\forall x:N\,\forall y:N\,(Sx=_NSy\Rightarrow x=_Ny)}{Arithmetic extension only; same NNO interpretation.}
\hbrule{rule:p10-internal-hol-ls1980-peano-induction}{Optional nonlogical (d), induction}{}{\vdash_{\varnothing}\forall u:PN\,((0\in u\wedge\forall x:N\,(x\in u\Rightarrow Sx\in u))\Rightarrow\forall y:N\,y\in u)}{Arithmetic extension only; same NNO interpretation. Bound \(u,x,y\) are distinct.}
\hbrule{rule:p10-internal-hol-ls1980-boolean}{Optional Boolean axiom}{}{\vdash_{\varnothing}\forall t:\Omega\,(t\vee\neg t)}{Classical extension only; the interpreting topos is Boolean. Independent of whether Peano rows are included.}
